\documentclass{amsart}
% For dealing with references we use the comment environment
\usepackage{verbatim}
\newenvironment{reference}{\comment}{\endcomment}
%\newenvironment{reference}{}{}
\newenvironment{slogan}{\comment}{\endcomment}
\newenvironment{history}{\comment}{\endcomment}

% For commutative diagrams we use Xy-pic
\usepackage[all]{xy}

% We use 2cell for 2-commutative diagrams.
\xyoption{2cell}
\UseAllTwocells

% We use multicol for the list of chapters between chapters
\usepackage{multicol}

% This is generally recommended for better output
\usepackage{lmodern}
\usepackage[T1]{fontenc}

% For cross-file-references

% Package for hypertext links:
\usepackage{hyperref}

% For any local file, say "hello.tex" you want to link to please

% Theorem environments.
%
\theoremstyle{plain}
\newtheorem{theorem}[subsection]{Theorem}
\newtheorem{proposition}[subsection]{Proposition}
\newtheorem{lemma}[subsection]{Lemma}

\theoremstyle{definition}
\newtheorem{definition}[subsection]{Definition}
\newtheorem{example}[subsection]{Example}
\newtheorem{exercise}[subsection]{Exercise}
\newtheorem{situation}[subsection]{Situation}

\theoremstyle{remark}
\newtheorem{remark}[subsection]{Remark}
\newtheorem{remarks}[subsection]{Remarks}

\numberwithin{equation}{subsection}

% Macros
%
\def\lim{\mathop{\mathrm{lim}}\nolimits}
\def\colim{\mathop{\mathrm{colim}}\nolimits}
\def\Spec{\mathop{\mathrm{Spec}}}
\def\Hom{\mathop{\mathrm{Hom}}\nolimits}
\def\Ext{\mathop{\mathrm{Ext}}\nolimits}
\def\SheafHom{\mathop{\mathcal{H}\!\mathit{om}}\nolimits}
\def\SheafExt{\mathop{\mathcal{E}\!\mathit{xt}}\nolimits}
\def\Sch{\mathit{Sch}}
\def\Mor{\mathop{\mathrm{Mor}}\nolimits}
\def\Ob{\mathop{\mathrm{Ob}}\nolimits}
\def\Sh{\mathop{\mathit{Sh}}\nolimits}
\def\NL{\mathop{N\!L}\nolimits}
\def\CH{\mathop{\mathrm{CH}}\nolimits}
\def\proetale{{pro\text{-}\acute{e}tale}}
\def\etale{{\acute{e}tale}}
\def\QCoh{\mathit{QCoh}}
\def\Ker{\mathop{\mathrm{Ker}}}
\def\Im{\mathop{\mathrm{Im}}}
\def\Coker{\mathop{\mathrm{Coker}}}
\def\Coim{\mathop{\mathrm{Coim}}}

% Boxtimes
%
\DeclareMathSymbol{\boxtimes}{\mathbin}{AMSa}{"02}

%
% Macros for moduli stacks/spaces
%
\def\QCohstack{\mathcal{QC}\!\mathit{oh}}
\def\Cohstack{\mathcal{C}\!\mathit{oh}}
\def\Spacesstack{\mathcal{S}\!\mathit{paces}}
\def\Quotfunctor{\mathrm{Quot}}
\def\Hilbfunctor{\mathrm{Hilb}}
\def\Curvesstack{\mathcal{C}\!\mathit{urves}}
\def\Polarizedstack{\mathcal{P}\!\mathit{olarized}}
\def\Complexesstack{\mathcal{C}\!\mathit{omplexes}}
% \Pic is the operator that assigns to X its picard group, usage \Pic(X)
% \Picardstack_{X/B} denotes the Picard stack of X over B
% \Picardfunctor_{X/B} denotes the Picard functor of X over B
\def\Pic{\mathop{\mathrm{Pic}}\nolimits}
\def\Picardstack{\mathcal{P}\!\mathit{ic}}
\def\Picardfunctor{\mathrm{Pic}}
\def\Deformationcategory{\mathcal{D}\!\mathit{ef}}

\setlength{\emergencystretch}{3em}
\begin{document}
\title[Stacks Project, Tag 04EB]{1629 --- Universal square-zero thickening of an arbitrary formally unramified ring map}
\maketitle
\noindent Copyright (C) 2005--2025 Johan de Jong. Stacks Project authors. Source: \href{https://stacks.math.columbia.edu/tag/04EB}{Tag 04EB}.

\noindent Editorial difficulty note (not author text). Difficulty H1: The proof constructs an initial square-zero thickening for every formally unramified ring map, without finite-generation assumptions, by splitting the infinite free polynomial cotangent module and correcting all polynomial-generator lifts uniquely through the inverse conormal differential. This explicit universal obstruction-killing construction goes beyond ordinary uniqueness of infinitesimal lifts..

\noindent Editorial provenance note (not author text). History: excerpt from revision \texttt{a0444\allowbreak 6e57e\allowbreak c1fbc\allowbreak 252a8\allowbreak 71afc\allowbreak ec775\allowbreak 2fb28\allowbreak 07b14}, algebra.tex, lines 40919--40997. Reference presentation was adapted; original-author context and core support blocks were retained or added. The mathematical wording is unchanged. Licensed under GNU FDL 1.2 or later; no invariant sections or cover texts. See ../licenses/COPYING. Context blocks reproduce author definitions and supporting results; they are not additional counted problems.

\begin{definition}
\label{definition-formally-unramified}
Let $R \to S$ be a ring map.
We say $S$ is {\it formally unramified over $R$} if for every
commutative solid diagram
$$
\xymatrix{
S \ar[r] \ar@{-->}[rd] & A/I \\
R \ar[r] \ar[u] & A \ar[u]
}
$$
where $I \subset A$ is an ideal of square zero, there exists
at most one dotted arrow making the diagram commute.
\end{definition}

\begin{definition}
\label{definition-derivation}
Let $\varphi : R \to S$ be a ring map and let $M$ be an $S$-module.
A {\it derivation}, or more precisely an
{\it $R$-derivation} into $M$ is a map $D : S \to M$
which is additive, annihilates elements of $\varphi(R)$,
and satisfies the {\it Leibniz rule}: $D(ab) = aD(b) + bD(a)$.
\end{definition}

\begin{definition}
\label{definition-differentials}
The pair $(\Omega_{S/R}, \text{d})$ is called the {\it module
of K\"ahler differentials} or the {\it module of differentials}
of $S$ over $R$.
\end{definition}

\noindent
Suppose that
\begin{equation}
\label{equation-functorial-omega}
\vcenter{
\xymatrix{
S \ar[r]_\varphi
&
S'
\\
R \ar[r]^\psi \ar[u]^\alpha
&
R' \ar[u]_\beta
}
}
\end{equation}
is a commutative diagram of rings. In this case there is a
natural map of modules of differentials fitting into the
commutative diagram
$$
\xymatrix{
\Omega_{S/R} \ar[r] &
\Omega_{S'/R'}
\\
S \ar[u]^{\text{d}} \ar[r]^{\varphi}
&
S' \ar[u]_{\text{d}}
}
$$
To construct the map just use the obvious map
between the presentations for $\Omega_{S/R}$ and $\Omega_{S'/R'}$.
Namely,
\begin{equation}
\label{equation-map-presentations}
\vcenter{
\xymatrix{
\bigoplus S'[(a', b')]
\oplus
\bigoplus S'[(f', g')]
\oplus
\bigoplus S'[r'] \ar[r]
&
\bigoplus S' [a'] \\
 \\
\bigoplus S[(a, b)]
\oplus
\bigoplus S[(f, g)]
\oplus
\bigoplus S[r] \ar[r]
\ar[uu]^{
\begin{matrix}
[(a, b)] \mapsto [(\varphi(a), \varphi(b))] \\
[(f, g)] \mapsto [(\varphi(f), \varphi(g))] \\
[r]\mapsto [\psi(r)]
\end{matrix}
} &
\bigoplus S[a] \ar[uu]_{[a] \mapsto [\varphi(a)]}
}
}
\end{equation}
The result is simply that $f\text{d}g \in \Omega_{S/R}$ is
mapped to $\varphi(f)\text{d}\varphi(g)$.

\begin{lemma}
\label{lemma-characterize-formally-unramified}
Let $R \to S$ be a ring map.
The following are equivalent:
\begin{enumerate}
\item $R \to S$ is formally unramified,
\item the module of differentials $\Omega_{S/R}$ is zero.
\end{enumerate}
\end{lemma}

\begin{proof}
Let $J = \Ker(S \otimes_R S \to S)$ be the kernel of
the multiplication map. Let $A_{univ} = S \otimes_R S/J^2$. Recall
that $I_{univ} = J/J^2$ is isomorphic to $\Omega_{S/R}$, see
Lemma \href{https://stacks.math.columbia.edu/tag/00RW}{lemma differentials diagonal (Tag 00RW)}. Moreover, the two $R$-algebra maps
$\sigma_1, \sigma_2 : S \to A_{univ}$, $\sigma_1(s) = s \otimes 1 \bmod J^2$,
and $\sigma_2(s) = 1 \otimes s \bmod J^2$ differ by the
universal derivation $\text{d} : S \to \Omega_{S/R} = I_{univ}$.

\medskip\noindent
Assume $R \to S$ formally unramified.
Then we see that $\sigma_1 = \sigma_2$.
Hence $\text{d}(s) = 0$ for all $s \in S$.
Hence $\Omega_{S/R} = 0$.

\medskip\noindent
Assume that $\Omega_{S/R} = 0$. Let $A, I, R \to A, S \to A/I$
be a solid diagram as in Definition \ref{definition-formally-unramified}.
Let $\tau_1, \tau_2 : S \to A$ be two dotted arrows making the
diagram commute. Consider the $R$-algebra map $A_{univ} \to A$
defined by the rule $s_1 \otimes s_2 \mapsto \tau_1(s_1)\tau_2(s_2)$.
We omit the verification that this is well defined. Since $A_{univ} \cong S$
as $I_{univ} = \Omega_{S/R} = 0$ we conclude that $\tau_1 = \tau_2$.
\end{proof}


\begin{lemma}
\label{lemma-differential-seq}
In diagram (\ref{equation-functorial-omega}),
suppose that $S \to S'$ is surjective with kernel $I \subset S$,
and assume that $R' = R$.
Then there is a canonical exact sequence of $S'$-modules
$$
I/I^2
\longrightarrow
\Omega_{S/R} \otimes_S S'
\longrightarrow
\Omega_{S'/R}
\longrightarrow
0
$$
The leftmost map is characterized by the rule that
$f \in I$ maps to $\text{d}f \otimes 1$.
\end{lemma}

\begin{proof}
The middle term is $\Omega_{S/R} \otimes_S S/I$.
For $f \in I$ denote $\overline{f}$ the image of $f$ in $I/I^2$.
To show that the map $\overline{f} \mapsto \text{d}f \otimes 1$
is well defined we just have to check that
$\text{d} f_1f_2 \otimes 1 = 0$ if $f_1, f_2 \in I$.
And this is clear from the Leibniz rule
$\text{d} f_1f_2 \otimes 1 = (f_1 \text{d}f_2 + f_2 \text{d} f_1 )\otimes 1 =
\text{d}f_2 \otimes f_1 + \text{d}f_1 \otimes f_2 = 0$.
A similar computation show this map is $S' = S/I$-linear.

\medskip\noindent
The map $\Omega_{S/R} \otimes_S S' \to \Omega_{S'/R}$
is the canonical $S'$-linear map associated to the
$S$-linear map $\Omega_{S/R} \to \Omega_{S'/R}$.
It is surjective because $\Omega_{S/R} \to \Omega_{S'/R}$
is surjective by Lemma \href{https://stacks.math.columbia.edu/tag/00RR}{lemma differential surjective (Tag 00RR)}.

\medskip\noindent
The composite of the two maps is zero because
$\text{d}f$ maps to zero in $\Omega_{S'/R}$
for $f \in I$. Note that exactness just says that
the kernel of $\Omega_{S/R} \to \Omega_{S'/R}$
is generated as an $S$-submodule by the submodule $I\Omega_{S/R}$ together
with the elements $\text{d}f$, with $f \in I$. We know by
Lemma \href{https://stacks.math.columbia.edu/tag/00RR}{lemma differential surjective (Tag 00RR)}
that this kernel is generated by the elements $\text{d}(a)$
where $\varphi(a) = \beta(r)$ for some $r \in R$.
But then $a = \alpha(r) + a - \alpha(r)$, so
$\text{d}(a) = \text{d}(a - \alpha(r))$. And
$a - \alpha(r) \in I$ since $\varphi(a - \alpha(r)) =
\varphi(a) - \varphi(\alpha(r)) = \beta(r) - \beta(r) = 0$.
We conclude the elements $\text{d}f$ with $f \in I$ already
generate the kernel as an $S$-module, as desired.
\end{proof}


\begin{lemma}
\label{lemma-universal-omega}
\begin{slogan}
Maps out of the module of differentials are the same as derivations.
\end{slogan}
The module of differentials of $S$ over $R$ has the following
universal property. The map
$$
\Hom_S(\Omega_{S/R}, M)
\longrightarrow
\text{Der}_R(S, M), \quad
\alpha
\longmapsto
\alpha \circ \text{d}
$$
is an isomorphism of functors.
\end{lemma}

\begin{proof}
By definition an $R$-derivation is a rule which associates
to each $a \in S$ an element $D(a) \in M$. Thus $D$ gives
rise to a map $[D] : \bigoplus S[a] \to M$. However, the conditions
of being an $R$-derivation exactly mean that $[D]$ annihilates
the image of the map in the displayed presentation of
$\Omega_{S/R}$ above.
\end{proof}



\medskip\noindent
Consider the following map of free $S$-modules
$$
\bigoplus\nolimits_{(a, b)\in S^2} S[(a, b)] \oplus
\bigoplus\nolimits_{(f, g)\in S^2} S[(f, g)] \oplus
\bigoplus\nolimits_{r\in R} S[r]
\longrightarrow
\bigoplus\nolimits_{a\in S} S[a]
$$
defined by the rules
$$
[(a, b)] \longmapsto [a + b] - [a] - [b],\quad
[(f, g)] \longmapsto [fg] -f[g] - g[f],\quad
[r]      \longmapsto [\varphi(r)]
$$
with obvious notation. Let $\Omega_{S/R}$ be the cokernel of this map.
There is a map $\text{d} : S \to \Omega_{S/R}$ which maps $a$ to the
class $\text{d}a$ of $[a]$ in the cokernel. This is an $R$-derivation
by the relations imposed on $\Omega_{S/R}$, in other words
$$
\text{d}(a + b) = \text{d}a + \text{d}b, \quad
\text{d}(fg) = f\text{d}g + g\text{d}f, \quad
\text{d}\varphi(r) = 0
$$
where $a,b,f,g \in S$ and $r \in R$.


\begin{lemma}
\label{lemma-universal-thickening}
Let $R \to S$ be a formally unramified ring map. There exists a surjection of
$R$-algebras $S' \to S$ whose kernel is an ideal of square zero with the
following universal property: Given any commutative diagram
$$
\xymatrix{
S \ar[r]_a & A/I \\
R \ar[r]^b \ar[u] & A \ar[u]
}
$$
where $I \subset A$ is an ideal of square zero, there is a unique $R$-algebra
map $a' : S' \to A$ such that $S' \to A \to A/I$ is equal to $S' \to S \to A/I$.
\end{lemma}

\begin{proof}
Choose a set of generators $z_i \in S$, $i \in I$ for $S$ as an $R$-algebra.
Let $P = R[\{x_i\}_{i \in I}]$ denote the polynomial ring on generators
$x_i$, $i \in I$. Consider the $R$-algebra map $P \to S$ which maps
$x_i$ to $z_i$. Let $J = \Ker(P \to S)$. Consider the map
$$
\text{d} : J/J^2 \longrightarrow \Omega_{P/R} \otimes_P S
$$
see
Lemma \ref{lemma-differential-seq}.
This is surjective since $\Omega_{S/R} = 0$ by assumption, see
Lemma \ref{lemma-characterize-formally-unramified}.
Note that $\Omega_{P/R}$ is free on $\text{d}x_i$, and hence the module
$\Omega_{P/R} \otimes_P S$ is free over $S$. Thus we may choose a splitting
of the surjection above and write
$$
J/J^2 = K \oplus \Omega_{P/R} \otimes_P S
$$
Let $J^2 \subset J' \subset J$ be the ideal of $P$ such that
$J'/J^2$ is the second summand in the decomposition above.
Set $S' = P/J'$. We obtain a short exact sequence
$$
0 \to J/J' \to S' \to S \to 0
$$
and we see that $J/J' \cong K$ is a square zero ideal in $S'$. Hence
$$
\xymatrix{
S \ar[r]_1 & S \\
R \ar[r] \ar[u] & S' \ar[u]
}
$$
is a diagram as above. In fact we claim that this is an initial object in
the category of diagrams. Namely, let $(I \subset A, a, b)$ be an arbitrary
diagram. We may choose an $R$-algebra map $\beta : P \to A$ such that
$$
\xymatrix{
S \ar[r]_1 & S \ar[r]_a & A/I \\
R \ar[r] \ar@/_/[rr]_b \ar[u] & P \ar[u] \ar[r]^\beta & A \ar[u]
}
$$
is commutative. Now it may not be the case that $\beta(J') = 0$, in other
words it may not be true that $\beta$ factors through $S' = P/J'$.
But what is clear is that $\beta(J') \subset I$ and
since $\beta(J) \subset I$ and $I^2 = 0$ we have $\beta(J^2) = 0$.
Thus the ``obstruction'' to finding a morphism from
$(J/J' \subset S', 1, R \to S')$ to $(I \subset A, a, b)$ is
the corresponding $S$-linear map $\overline{\beta} : J'/J^2 \to I$.
The choice in picking $\beta$ lies in the choice of $\beta(x_i)$.
A different choice of $\beta$, say $\beta'$, is gotten by taking
$\beta'(x_i) = \beta(x_i) + \delta_i$ with $\delta_i \in I$.
In this case, for $g \in J'$, we obtain
$$
\beta'(g) =
\beta(g) + \sum\nolimits_i \delta_i \beta(\frac{\partial g}{\partial x_i})
$$
Since the map $\text{d}|_{J'/J^2} : J'/J^2 \to \Omega_{P/R} \otimes_P S$
given by $g \mapsto \frac{\partial g}{\partial x_i}\text{d}x_i \otimes 1$
is an isomorphism by construction, we see that there is a unique choice
of $\delta_i \in I$ such that $\beta'(g) = 0$ for all $g \in J'$.
(Namely, $\delta_i$ is $-\overline{\beta}(g)$ where $g \in J'/J^2$
is the unique element mapped to $\text{d}x_i \otimes 1$
by $\text{d}|_{J'/J^2}$.) The uniqueness of the solution implies the
uniqueness required in the lemma.
\end{proof}
\end{document}
